\documentclass[11pt]{article}
\usepackage[margin=1in]{geometry}
\usepackage{amsmath,amsthm,amssymb,mathtools}
\usepackage[margin=1in]{geometry}  % already defined in main document
\usepackage{amsmath,amsthm,amssymb,mathtools}  % already defined

\usepackage{graphicx}
\usepackage{xcolor}
\usepackage{listings}
\usepackage{float}


\title{Calvo Pricing in the Simple Monetary Sticky-Price Model}
\author{}
\date{}

\begin{document}
\maketitle


\section{Setup}

We keep the \emph{real side} of the lecture notes unchanged and only replace the pricing friction.
The original notes impose a simple environment with no capital in production, linear technology,
and a money-demand relation that implies consumption equals real balances in logs. In particular,
we retain:
\begin{align}
y_t &= c_t, \\
m_t - p_t &= c_t, \\
w_t &= \gamma y_t.
\end{align}
Combining the first two equations gives
\begin{equation}
y_t = m_t - p_t,
\end{equation}
and hence
\begin{equation}
w_t = \gamma (m_t - p_t).
\end{equation}

Here:
\begin{itemize}
    \item $m_t$ is the log money stock,
    \item $p_t$ is the log aggregate price level,
    \item $y_t$ is log output,
    \item $w_t$ is the log real wage,
    \item $\gamma>0$ measures how strongly real marginal cost responds to output.
\end{itemize}

The only change relative to the lecture notes is that we replace the fixed-duration pricing scheme
($N=2$ staggered pricing) with \emph{Calvo pricing}.

\section{Calvo pricing}

In each period, a firm receives an opportunity to reset its price with probability $1-\theta$.
With probability $\theta$, it cannot reoptimize and must keep its previous nominal price.

Let $p_t^*$ denote the optimal reset price chosen by firms that can adjust in period $t$.
Under Calvo pricing, a firm that resets in period $t$ chooses $p_t^*$ to maximize expected discounted
profits over the random duration for which the price remains in effect.

Because we keep the rest of the model unchanged, the firm's nominal marginal cost in logs is
\begin{equation}
p_t + w_t.
\end{equation}
Therefore, the linearized optimal reset-price equation is
\begin{equation}
p_t^* = (1-\beta\theta)\sum_{j=0}^{\infty} (\beta\theta)^j E_t\bigl[p_{t+j}+w_{t+j}\bigr].
\label{eq:reset_infinite}
\end{equation}
This is the standard Calvo reset-price formula: the newly chosen price is a discounted weighted
average of current and future nominal marginal costs.

Equation \eqref{eq:reset_infinite} can be written recursively as
\begin{equation}
p_t^* = (1-\beta\theta)(p_t+w_t) + \beta\theta E_t p_{t+1}^*.
\label{eq:reset_recursive}
\end{equation}

\section{Aggregate price dynamics}

Under Calvo pricing, the aggregate price level is a weighted average of firms that do not adjust
and firms that do adjust. In log-linear form,
\begin{equation}
p_t = \theta p_{t-1} + (1-\theta)p_t^*.
\label{eq:aggregate_price}
\end{equation}
Define inflation by
\begin{equation}
\pi_t \equiv p_t - p_{t-1}.
\end{equation}
Then \eqref{eq:aggregate_price} implies
\begin{equation}
\pi_t = (1-\theta)(p_t^* - p_{t-1}).
\label{eq:inflation_reset_gap}
\end{equation}
From this,
\begin{align}
p_t^* - p_t
&= p_t^* - \bigl[\theta p_{t-1} + (1-\theta)p_t^*\bigr] \\
&= \theta(p_t^* - p_{t-1}) \\
&= \frac{\theta}{1-\theta}\pi_t.
\end{align}
Thus,
\begin{equation}
p_t^* - p_t = \frac{\theta}{1-\theta}\pi_t.
\label{eq:gap_pi}
\end{equation}

\section{Derivation of the Phillips curve}

Subtract $p_t$ from both sides of \eqref{eq:reset_recursive}:
\begin{align}
p_t^* - p_t
&= (1-\beta\theta)w_t + \beta\theta E_t p_{t+1}^* - \beta\theta p_t \\
&= (1-\beta\theta)w_t + \beta\theta E_t\bigl[(p_{t+1}^*-p_{t+1}) + (p_{t+1}-p_t)\bigr] \\
&= (1-\beta\theta)w_t + \beta\theta E_t(p_{t+1}^* - p_{t+1}) + \beta\theta E_t\pi_{t+1}.
\end{align}
Now use \eqref{eq:gap_pi} both at time $t$ and $t+1$:
\begin{equation}
p_t^* - p_t = \frac{\theta}{1-\theta}\pi_t,
\qquad
p_{t+1}^* - p_{t+1} = \frac{\theta}{1-\theta}\pi_{t+1}.
\end{equation}
Substituting these into the previous equation yields
\begin{equation}
\frac{\theta}{1-\theta}\pi_t
=
(1-\beta\theta)w_t
+
\beta\theta E_t\left(\frac{\theta}{1-\theta}\pi_{t+1}\right)
+
\beta\theta E_t\pi_{t+1}.
\end{equation}
Multiply both sides by $(1-\theta)/\theta$:
\begin{equation}
\pi_t
=
\frac{(1-\theta)(1-\beta\theta)}{\theta}w_t
+
\beta E_t\pi_{t+1}.
\end{equation}
Define
\begin{equation}
\kappa \equiv \frac{(1-\theta)(1-\beta\theta)}{\theta}.
\end{equation}
Then the Phillips curve is
\begin{equation}
\pi_t = \beta E_t\pi_{t+1} + \kappa w_t.
\label{eq:nkpc_w}
\end{equation}
Using $w_t=\gamma y_t$, we obtain
\begin{equation}
\boxed{\pi_t = \beta E_t\pi_{t+1} + \kappa\gamma y_t.}
\label{eq:nkpc_y}
\end{equation}

This is the Calvo analogue of the pricing equation in the lecture notes. The key difference is that
the fixed $N=2$ pricing block is replaced by a forward-looking Phillips curve.

\section{Combining the Phillips curve with money demand}

Since
\begin{equation}
y_t = m_t - p_t,
\end{equation}
we may write
\begin{equation}
p_t = m_t - y_t.
\end{equation}
Therefore inflation is
\begin{align}
\pi_t
&= p_t - p_{t-1} \\
&= (m_t - y_t) - (m_{t-1} - y_{t-1}) \\
&= \Delta m_t - (y_t - y_{t-1}),
\end{align}
where
\begin{equation}
\Delta m_t \equiv m_t - m_{t-1}.
\end{equation}
Hence
\begin{equation}
\boxed{\pi_t = y_{t-1} + \Delta m_t - y_t.}
\label{eq:pi_output}
\end{equation}

Substitute \eqref{eq:pi_output} into \eqref{eq:nkpc_y}:
\begin{equation}
y_{t-1} + \Delta m_t - y_t
=
\beta E_t\pi_{t+1} + \kappa\gamma y_t.
\end{equation}
But from \eqref{eq:pi_output} one period ahead,
\begin{equation}
\pi_{t+1} = y_t + \Delta m_{t+1} - y_{t+1}.
\end{equation}
Thus,
\begin{equation}
y_{t-1} + \Delta m_t - y_t
=
\beta E_t\bigl[y_t + \Delta m_{t+1} - y_{t+1}\bigr]
+
\kappa\gamma y_t.
\label{eq:master}
\end{equation}

Equation \eqref{eq:master} is the equilibrium law of motion for output in the Calvo version of the model.

\section{Special case: money is a random walk}

Suppose the money stock follows
\begin{equation}
m_t = m_{t-1} + \varepsilon_t,
\qquad
E_t\varepsilon_{t+1}=0.
\end{equation}
Then
\begin{equation}
\Delta m_t = \varepsilon_t,
\qquad
E_t\Delta m_{t+1}=0.
\end{equation}
Substituting this into \eqref{eq:master} gives
\begin{equation}
y_{t-1} + \varepsilon_t - y_t
=
\beta\bigl(y_t - E_t y_{t+1}\bigr)
+
\kappa\gamma y_t.
\label{eq:rw_equation}
\end{equation}

Now conjecture a linear solution of the form
\begin{equation}
y_t = a y_{t-1} + b\varepsilon_t.
\label{eq:guess}
\end{equation}
Then
\begin{equation}
E_t y_{t+1} = a y_t.
\end{equation}
Substituting into \eqref{eq:rw_equation},
\begin{equation}
y_{t-1} + \varepsilon_t - y_t
=
\beta(1-a)y_t + \kappa\gamma y_t.
\end{equation}
Rearrange:
\begin{equation}
y_{t-1} + \varepsilon_t
=
\bigl[1 + \beta(1-a) + \kappa\gamma\bigr]y_t.
\end{equation}
Using the guess \eqref{eq:guess},
\begin{equation}
y_{t-1} + \varepsilon_t
=
\bigl[1 + \beta(1-a) + \kappa\gamma\bigr](a y_{t-1} + b\varepsilon_t).
\end{equation}
Matching coefficients on $y_{t-1}$ and $\varepsilon_t$ implies
\begin{equation}
a=b,
\end{equation}
and
\begin{equation}
1 = \bigl[1 + \beta(1-a) + \kappa\gamma\bigr]a.
\end{equation}
Expanding,
\begin{equation}
1 = a + \beta a - \beta a^2 + \kappa\gamma a,
\end{equation}
or equivalently
\begin{equation}
\boxed{\beta a^2 - (1+\beta+\kappa\gamma)a + 1 = 0.}
\label{eq:a_quadratic}
\end{equation}
The stable root is the one inside the unit circle:
\begin{equation}
\boxed{
a
=
\frac{(1+\beta+\kappa\gamma)
-
\sqrt{(1+\beta+\kappa\gamma)^2 - 4\beta}}
{2\beta}.
}
\label{eq:a_root}
\end{equation}
Therefore output obeys
\begin{equation}
\boxed{y_t = a y_{t-1} + a\varepsilon_t.}
\label{eq:y_solution}
\end{equation}
Using \eqref{eq:pi_output},
\begin{equation}
\pi_t = y_{t-1} + \varepsilon_t - y_t
= (1-a)(y_{t-1} + \varepsilon_t).
\end{equation}
Hence
\begin{equation}
\boxed{\pi_t = (1-a)(y_{t-1} + \varepsilon_t).}
\label{eq:pi_solution}
\end{equation}

\section{Interpretation}

Compared with the original fixed-duration pricing model in the lecture notes, the Calvo version
changes the dynamic structure of inflation and output.

In the lecture notes, the fixed $N=2$ pricing scheme implies a scalar root
\begin{equation}
a_{\text{Taylor}} = \frac{1-\sqrt{\gamma}}{1+\sqrt{\gamma}},
\end{equation}
which becomes negative when $\gamma>1$. That negative root is what generates the undesirable
oscillatory or non-persistent behavior emphasized in the notes.

Under Calvo pricing instead, the persistence parameter $a$ is determined by
\begin{equation}
\beta a^2 - (1+\beta+\kappa\gamma)a + 1 = 0.
\end{equation}
The stable root satisfies
\begin{equation}
0<a<1.
\end{equation}
Therefore a monetary expansion $\varepsilon_t>0$ raises output on impact:
\begin{equation}
\frac{\partial y_t}{\partial \varepsilon_t} = a > 0.
\end{equation}
So replacing fixed-duration pricing by Calvo pricing eliminates the negative-root pathology.

However, the model still retains the same underlying force emphasized in the lecture notes:
if $\gamma$ is very large, then the term $\kappa\gamma$ is large, and the stable root $a$ becomes small.
That means the positive effect of monetary shocks on output is not very persistent.

Thus, changing only the nominal rigidity from fixed-duration pricing to Calvo pricing improves the
qualitative behavior of the model, but it does not remove the deeper issue that real marginal cost
responds too strongly to output.

\end{document}